% Chapter 7, Section 6

\section{Other Regularization Techniques}
\label{sec:other-regularization}

\subsection{Intuition: Many Small Guards Against Overfitting}

Beyond penalties and normalization, there are practical techniques that act like small guards during training. Each introduces a mild constraint or noise that nudges the model away from brittle solutions and encourages smoother decision boundaries.

\subsection{Label Smoothing}

Replace hard targets with smoothed distributions:
\begin{equation}
y'_k = (1 - \epsilon) y_k + \frac{\epsilon}{K}
\end{equation}

Prevents overconfident predictions.

\subsection{Gradient Clipping}

Limit gradient magnitude to prevent exploding gradients:

\textbf{Clipping by value:}
\begin{equation}
g \leftarrow \max(\min(g, \theta), -\theta)
\end{equation}

\textbf{Clipping by norm:}
\begin{equation}
g \leftarrow \frac{g}{\max(1, \|g\| / \theta)}
\end{equation}

\subsection{Stochastic Depth}

Randomly skip layers during training (for very deep networks).

\subsection{Mixup}

Train on convex combinations of examples:
\begin{align}
\tilde{\vect{x}} &= \lambda \vect{x}_i + (1-\lambda) \vect{x}_j \\
\tilde{y} &= \lambda y_i + (1-\lambda) y_j
\end{align}

where $\lambda \sim \text{Beta}(\alpha, \alpha)$.

\subsection{Adversarial Training}

Add adversarially perturbed examples to training:
\begin{equation}
\vect{x}_{\text{adv}} = \vect{x} + \epsilon \cdot \text{sign}(\nabla_{\vect{x}} L(\vect{x}, y))
\end{equation}

% Index entries
\index{label smoothing}
\index{gradient clipping}
\index{stochastic depth}
\index{mixup}
\index{adversarial training}
